% !PRINTED 209
\weberpage{209}

Dass auch umgekehrt jeder solche Ausdruck, wenn er in
den Co\"efficienten $a$ homogen ist, eine ganzzahlige Invariante
darstellt, ist von selbst klar\footnote{Die Theorie der Invarianten findet man ausf\"uhrlich dargestellt in
den Werken: \gesperrt{Clebsch}, ,,Theorie der bin\"aren algebraischen Formen``.
Leipzig 1872. \gesperrt{Fa\`a di Bruno}, ,,Einleitung in die Theorie der bin\"aren
Formen``. Deutsch von \gesperrt{Walter}. Leipzig 1881. \gesperrt{P. Gordans}, ,,Vorlesungen
\"uber Invariantentheorie``, herausgegeben von \gesperrt{Kerschensteiner}. Leipzig
1885. Vgl. auch \gesperrt{Franz Meyer}, ,,Bericht \"uber den gegenw\"artigen Stand
der Invariantentheorie`` im Jahresbericht der deutschen Mathematiker-
Vereinigung 1890/91 (Berlin 1892).}.

% Printer's furniture excluded by policy: footer "Weber, Algebra. I." and gathering number 14.
